\section{The Cassels--Tate comparison}
\label{sec:chi-ct-p}

The first supporting comparison used by the closure is the
Cassels--Tate Pfaffian comparison~\citep{Cassels1962,PoonenStoll1999}.  The import stack introduced in
Definition~\ref{def:closure:chi-ct-p-import} packages the relevant external
results as hypotheses: the Selmer-complex inputs of Sakamoto~\citep{Sakamoto2018}
and Macias Castillo--Sano~\citep{MaciasCastilloSano2026}, the determinant
and orientation inputs of Knudsen--Mumford and Nekov\'a\v{r}, and Flach's
pairing comparison.  Fisher--Schaefer--Stoll~\citep{FisherSchaeferStoll2010}
fix the classical 2-Selmer pairing normalization; they are not used as a
Fitting-ideal theorem.  This section states the
comparison in the form consumed by the closure.  It is conditional on
that stack; no part of the argument below reproves the imported
arithmetic geometry.

\begin{theorem}[$\chi_{CT,p}$ comparison (v5)]\label{thm:closure:chi-ct-p-comparison}
Let \(E/\Q\) be an elliptic curve such that \(\Sha(E)[p^\infty]\) is
finite, the Cassels--Tate pairing is nondegenerate, the standard
Sakamoto and Macias Castillo--Sano Selmer-complex hypotheses hold, and
\(p\nmid \Tam(E/\Q)\).  Conditional on the imported \(\chiCTp\)
structure, there is a canonical orientation-respecting Pfaffian
comparison map
\[
  \Pf_{\mathrm{Nek},p}^{\mathrm{or},\sqrt{\ }}:
  \sqrt{\Fitt}_{\mathrm{CT}}
    \bigl(H^1(C_p(E))/\mathrm{div},\,U_{2,2}\bigr)
  \longrightarrow
  \mathrm{vctp}_p
\]
from the square-root Cassels--Tate Fitting ideal on the quotient by the
divisible subgroup, paired by the Nekov\'a\v{r} cup-product
\(U_{2,2}\) identified with the Flach pairing, to the
Cassels--Tate Pfaffian scalar line \(\mathrm{vctp}_p\).  This map sends
the square-root Stark generator~\citep{MazurRubin2004} \(\sqrt{F}_E^p\) to the
Cassels--Tate height \(h_p^{\mathrm{CT}}(E)\), and the Nekov\'a\v{r}
sign convention satisfies \(\Sigma_{\mathrm{NekCT}}=+1\).\footnote{Tracked in the formalization harness as a conditional schema
\texttt{SixBirdsBSD.\allowbreak Closure.\allowbreak ChiCTp.\allowbreak chiCTpComparison}
-- a checked projection over an explicit proof-carrying record
(conditional on the imported \(\chiCTp\) stack), not an independent Lean
derivation. See App~\ref{app:formalization}.}
\end{theorem}

\begin{proof}
Assume the imported \(\chiCTp\) structure.  Its named components supply
the ingredients in the statement.  The Selmer-complex hypotheses and
the Cassels--Tate nondegeneracy input place the quotient
\(H^1(C_p(E))/\mathrm{div}\) in the self-dual setting in which a
Pfaffian comparison is defined.  The determinant-line and orientation
inputs identify the square-root Fitting ideal
\(\sqrt{\Fitt}_{\mathrm{CT}}(H^1(C_p(E))/\mathrm{div},U_{2,2})\) as the
source of the comparison map.  Here the square-root Fitting ideal is
the square-root determinant object attached to the Cassels--Tate
pairing; the notation \(U_{2,2}\) records the Nekov\'a\v{r}
cup-product, and the imported Flach comparison identifies it with the
classical Cassels--Tate pairing used on the target side.

The Pfaffian input in the stack then gives the canonical
orientation-respecting map
\[
  \Pf_{\mathrm{Nek},p}^{\mathrm{or},\sqrt{\ }}:
  \sqrt{\Fitt}_{\mathrm{CT}}
    \bigl(H^1(C_p(E))/\mathrm{div},U_{2,2}\bigr)
  \longrightarrow
  \mathrm{vctp}_p .
\]
The square-root Stark-generator identification supplies the image
formula: the distinguished generator \(\sqrt{F}_E^p\) maps to the
Cassels--Tate height \(h_p^{\mathrm{CT}}(E)\).  Finally, the sign
closure component of the same imported record fixes the orientation
sign as \(\Sigma_{\mathrm{NekCT}}=+1\).  Combining these imported
components proves exactly the displayed comparison.  The proof is
therefore an assembly under the stated hypotheses; it is precisely as
strong as the imported theorem stack and does not rederive the
Sakamoto, Macias Castillo--Sano, Nekov\'a\v{r}, Flach, or classical
Cassels--Tate comparison inputs.
\end{proof}

% ===========================================================================
